\section*{Purpose and Character of These Notes}
\addcontentsline{toc}{section}{Purpose and Character of These Notes}

Classical field theory is often presented in an extremely compressed form. A few symbols are introduced, an action is written down, a variation is performed in one line, and a field equation appears. To an experienced reader, such a presentation may be efficient. To a student meeting the subject for the first time, it can make the theory look like a collection of formulas that arrive without a clear reason, operate by hidden rules, and disappear before their meaning has been examined.

These notes are designed to avoid that style of exposition. Their purpose is not merely to collect standard formulas from analytical mechanics, special relativity, electrodynamics, and gauge theory. Their purpose is to show how the main constructions of classical field theory are built, why they are introduced, what assumptions they use, how the calculations are carried out, and how the final expressions are interpreted physically and mathematically.

A formula will therefore not be treated as the end of an explanation. Whenever possible, it will be presented as one stage in a sequence:

\begin{enumerate}
    \item identify the physical or mathematical problem;
    \item determine what objects must be introduced;
    \item state the assumptions and conventions;
    \item derive the relevant relation;
    \item test it on a simple model;
    \item interpret the result;
    \item examine limiting cases and possible failures.
\end{enumerate}

The central principle of the course is that a student should not only recognize an equation, but should know what to do with it. For example, it is not enough to state the Euler--Lagrange field equations. A student should see how the variation of a field is defined, how derivatives of the variation arise, where integration by parts is used, why a boundary term disappears, how indices are managed, and how the resulting equation is applied to an explicit Lagrangian density. Similarly, it is not enough to write down Noether's theorem. The theorem must be used to construct an actual current, verify its divergence, and identify the associated conserved charge.

The same standard will be applied to electrodynamics. The electromagnetic field tensor will not be introduced as an isolated matrix to memorize. It will be built from the four-potential, its components will be related explicitly to the electric and magnetic fields, its transformation properties will be examined, and its invariants will be calculated. Maxwell's equations will then be derived from an action, rather than merely restated in covariant notation.

These notes are meant to serve two related purposes. During the lecture, they provide a structured path through the material and a reliable source of notation, derivations, and examples. After the lecture, they should function as a readable reconstruction of the argument. A student who did not follow one step on the board should be able to return to the notes and see how that step fits into the complete calculation.

The exposition will therefore contain several recurring elements.

\subsection*{Motivation before formalism}

A new object will be introduced because it solves a problem or expresses a principle. The field concept will be motivated by locality and continuous degrees of freedom. The action will be motivated as a compact way to encode dynamics. Four-dimensional notation will be motivated by Lorentz symmetry. Gauge freedom will be motivated by the non-uniqueness of the potentials that describe the same physical electromagnetic field.

\subsection*{Definitions connected to examples}

Definitions will be followed by concrete cases. Scalar fields will be represented by temperature, density, or a real dynamical field. Vector fields will be connected with velocity fields and electric fields. Tensor fields will be introduced through objects that are actually used later, rather than through formal classification alone.

\subsection*{Complete derivations}

Important formulas will be derived in sufficient detail that the logical structure is visible. This does not mean that every elementary algebraic step must be expanded indefinitely. It means that no essential conceptual step should be hidden behind phrases such as ``it is easy to show'' or ``after some manipulation.'' When a sign, boundary condition, index contraction, or change of variables matters, it will be displayed.

\subsection*{Simple proofs and checks}

Not every statement in physics requires a theorem-proof format, but important claims should be justified. We will prove elementary identities when they clarify the structure of the theory, verify conservation laws directly, check gauge invariance explicitly, and compare covariant formulas with familiar three-dimensional expressions. Dimensional analysis, symmetry checks, and limiting cases will be used as routine diagnostic tools.

\subsection*{Worked calculations}

The notes will include explicit calculations rather than only general rules. A free scalar field will be solved by a plane-wave ansatz. A Noether current will be computed for a specified transformation. The electromagnetic field tensor will be written component by component. The Green function for the Poisson equation will be applied to a point source. The radiation field of an oscillating dipole will be developed far enough to calculate the angular distribution and total emitted power.

\subsection*{Physical interpretation}

After a derivation, we will ask what the result means. What quantity is conserved? What propagates? Which variables are measurable? Which variables contain redundancy? What changes under a transformation, and what remains invariant? A correct calculation without interpretation is incomplete, because field theory is not only a formal language but also a way of organizing physical information.

\subsection*{Connections between chapters}

The course is designed as a chain of constructions. The action principle for particles becomes the action principle for fields. Translation symmetry becomes energy-momentum conservation. The scalar field provides a simple model before the electromagnetic field is introduced. The electromagnetic potential becomes the first example of a gauge connection. The field strength becomes the first example of curvature. These connections will be stated explicitly so that the course does not appear as twelve unrelated topics.

The level of the notes is therefore deliberately intermediate between a concise reference text and a fully expanded textbook. They will remain mathematically precise, but they will also explain why the mathematics is being used. They will contain derivations and examples, but they will avoid turning every section into a long catalogue of exercises. The objective is a coherent set of lecture notes that can be read, checked, and used.
